\begin{tikzpicture}\begin{axis}[otLivreGraphe,width=100mm,height=51mm,
axis lines=left,xmin=0,xmax=8,ymin=0,ymax=7,
xtick={0,2,4,6,8},ytick={0,2,4,6},xlabel={$t\ (\mathrm h)$},ylabel={$C\ (\mathrm{mg/L})$},
every axis x label/.style={at={(axis description cs:1,0)},anchor=north east,yshift=-4mm,font=\small,text=otMuted},
every axis y label/.style={at={(axis description cs:0,1)},anchor=south west,yshift=1mm,font=\small,text=otMuted}]
\addplot[otNavy,thick,domain=0:8,samples=90]{8*x*exp(-.5*x)};
\addplot[only marks,otGreenDark]coordinates{(2,5.886071)};\draw[otGreenDark,dashed](axis cs:2,0)--(axis cs:2,5.886071);
\end{axis}\end{tikzpicture}
\par\smallskip\legende{Le changement de signe de la dérivée prouve le maximum observé à $t=2$ h.}
